\documentclass[11pt, a4paper]{article}

\usepackage{lmodern, amsmath, amssymb, amsthm, mathtools, mathrsfs, enumitem, hyperref, cleveref, geometry, microtype, suffix, graphicx, float, caption, subcaption, booktabs, multirow, siunitx, setspace, fancyhdr, cite, tikz, pgfplots, titlesec}
\usepackage[most]{tcolorbox}
\usepackage[utf8]{inputenc}
\usepackage[T1]{fontenc}

\pgfplotsset{compat=1.18}
\usetikzlibrary{3d, arrows.meta}

\newtheorem{definition}{Definition}
\newtheorem{proposition}{Proposition}
\newtheorem{remark}{Remark}
\newtheorem{example}{Example}
\newtheorem{theorem}{Theorem}
\newtcolorbox{boxdef}{enhanced}
\newtcolorbox{boxthem}{enhanced}
\newtcolorbox{boxprop}{enhanced}
\newtcolorbox{boxexa}{enhanced}

\geometry{margin=1.7cm}
\pagestyle{fancy}
\fancyhf{}
\fancyhead[L]{$\mathfrak{Analysis}$}
\fancyhead[C]{$\mathfrak{\&}$}
\fancyhead[R]{$\mathfrak{Mathematics}$}
\fancyfoot[C]{\thepage}
\renewcommand{\headrulewidth}{0.3pt}
\renewcommand{\footrulewidth}{0pt}

\newcommand{\N}{\mathbb{N}}
\newcommand{\R}{\mathbb{R}}
\newcommand{\Z}{\mathbb{Z}}
\newcommand{\Q}{\mathbb{Q}}
\newcommand{\C}{\mathbb{C}}
\newcommand{\E}{\mathbb{E}}
\newcommand{\F}{\mathbb{F}}
\newcommand{\pd}[2]{\dfrac{\partial #1}{\partial #2}}
\newcommand{\pdd}[3]{\dfrac{\partial^2 #1}{\partial #2\,\partial #3}}
\newcommand{\pddx}[2]{\dfrac{\partial^2 #1}{\partial #2^2}}
\newcommand{\dd}[2]{\dfrac{\mathrm{d} #1}{\mathrm{d} #2}}
\newcommand{\ddd}[2]{\dfrac{\mathrm{d}^2 #1}{\mathrm{d} #2^2}}
\newcommand{\diff}{\,\mathrm{d}}
\newcommand{\vect}[1]{\begin{pmatrix}#1\end{pmatrix}}
\newcommand{\abs}[1]{\lvert#1\rvert}
\newcommand{\set}[1]{\{#1\}}
\newcommand{\suchthat}{\;:\;}

\title{$\mathfrak{Analysis\ in\ Mathematics}$}
\author{}
\date{}

\begin{document}

\vspace*{2cm}
\begin{center}
    \huge{$\mathfrak{Analysis\ in\ Mathematics}$}
    \vspace*{3cm}
    \tableofcontents
    \newpage
\end{center}

\section{Integrals}
\documentclass[12pt]{article}
\usepackage[utf8]{inputenc}
\usepackage[T1]{fontenc}
\usepackage{amsmath, amssymb, amsthm}
\usepackage{siunitx}
\usepackage{tikz}
\usepackage{pgfplots}
\usepackage{geometry}
\usepackage{titlesec}
\geometry{a4paper, margin=1in}

\pgfplotsset{compat=1.18}
\usetikzlibrary{3d, arrows.meta}

% Theorem environments
\newtheorem{definition}{Definition}
\newtheorem{proposition}{Proposition}
\newtheorem{remark}{Remark}
\newtheorem{example}{Example}
\newtheorem{theorem}{Theorem}

% Paragraph formatting
\setlength{\parindent}{1.5em}
\setlength{\parskip}{1em}

% Add extra vertical space
\titleformat{\section}{\normalfont\Large\bfseries}{\thesection}{1em}{}
\titlespacing*{\section}{0pt}{2ex plus 1ex minus .2ex}{1ex plus .2ex}

\title{Chapter 01: Double integrals }
\author{Notes from prof zeghlaoui course}

\hbadness=10000
\hfuzz=\maxdimen

\newcommand{\N}{\mathbb{N}}
\newcommand{\R}{\mathbb{R}}
\newcommand{\Z}{\mathbb{Z}}
\newcommand{\Q}{\mathbb{Q}}
\newcommand{\C}{\mathbb{C}}
\newcommand{\E}{\mathbb{E}}
\newcommand{\F}{\mathbb{F}}

% --- Derivatives ---
\newcommand{\pd}[2]{\dfrac{\partial #1}{\partial #2}}
\newcommand{\pdd}[3]{\dfrac{\partial^2 #1}{\partial #2\,\partial #3}}
\newcommand{\pddx}[2]{\dfrac{\partial^2 #1}{\partial #2^2}}
\newcommand{\dd}[2]{\dfrac{\mathrm{d} #1}{\mathrm{d} #2}}
\newcommand{\ddd}[2]{\dfrac{\mathrm{d}^2 #1}{\mathrm{d} #2^2}}
\newcommand{\diff}{\,\mathrm{d}}

%% vertical vector macro
\newcommand{\vect}[1]{\begin{pmatrix}#1\end{pmatrix}}


\begin{document}

\maketitle

\vspace{2em}
\section{Reminder:}
Recall that if $f: \left[a,b\right] \to \R$ is bounded.

for any subdivision $\sigma = \{ x_0=a <x_1< x_2<...< x_{n+1} = b\}$ of  $\left[a,b\right]$

\[S(f,\sigma)= \sum_{i = 1}^{n}M_i(x_i - x_{i-1}) \quad \text{and} \quad s(f,\sigma)= \sum_{i = 1}^{n}m_i(x_i - x_{i-1})\]
\[\text{where } m_i = \inf_{x \in [x_{i-1},x_i]} f(x), \quad M_i = \sup_{x \in [x_{i-1},x_i]} f(x)\]

called respectively, the upper and the lower Darboux sums of $f$
on $\left[a, b\right]$ with respect to the subdivision $\sigma$;

$S$: the area of the region plane delimited by:
\[
\begin{cases}
  x=a \quad y = 0 \\
  x=b \quad y = f(x) \\
\end{cases}
\]

(graph showing the lower and upper of an integral)

\[s(f,\sigma) \le S \le S(f,\sigma)\quad \text{if } \sigma_1 \subset \sigma_2,\text{ then:}\]
\[s(f,\sigma_1)  \le s(f,\sigma_2) \le S \le S(f,\sigma_2) \le S(f,\sigma_1)\]

\[
S^+ = \inf (S(f,\sigma))\quad , \quad S^- = \sup (s(f,\sigma))
\]

\[
S^- \le S \le S^+
\]

\begin{definition}
  We say that $f$ is Riemann integrable on $\left[a,b\right]$ if :
  \[S^+ = S^- = S = \int_a^b f(x)dx\]
\end{definition}

\begin{definition}
  Riemann sums

  Let $\sigma$ a subdivision of $\left[a,b\right]$ and $\zeta = \{\zeta_1,\dots,\zeta_n\} \quad / \zeta_i \in \left[x_{i-1}, x_i\right]$

  \[s(f,\sigma) \le R(f,\sigma,\zeta) = \sum_{i=1}^{n}f(\zeta_i)(x_i-x_{i-1}) \le S(f,\sigma)\]

  If $f: \left[a,b\right] \to \R$ is Riemann integrable, then:

  $\exists (\sigma_n)$ a sequence of subdivision of $\left[a,b\right]$

$\exists (\zeta_n)$ associated with 
$(\sigma_n)[a,b]\quad \zeta_n = (\zeta_n^i),\ \zeta_n^i \in [x_{i-1}^n, x_i^n]$

  $\lim\limits_{n \to \infty}R(f, \sigma_n,\zeta_n) = \int_a^bf(x)dx$

\underline{In particular:}
\[\sigma_n = \{x_i^n= a + \frac{i}{n}(b-a): i\in {0,\dots,n}\}\]
\[\lim\limits_{n \to \infty} \frac{b-a}{n} \sum_{i = 1}^{n} f(\zeta_i) = \int_a^b f(x)dx\]
\[\lim\limits_{n \to \infty} \frac{b-a}{n} \sum_{i = 1}^{n}f(a+ \frac{i}{n}(b-a)) = \int_a^bf(x)dx \quad, x_i^n-x^n_{i-1} = \frac{b-a}{n}\]
\end{definition}

\section{Double integrals:}
\subsection{Definition and first properties:}

\begin{definition}
  Darboux sums- $\quad $ Riemann sums:

  Let $f: \Delta := \left[a,b\right]\times\left[c,d\right] \to \R $  be a bounded function

  1) A subdivision $\sigma$ of $\Delta$ is a set $\sigma = \{\Delta_{ij} = \left[x_{i-1}, x_i\right]\times\left[y_{j-1}, y_j\right] \}$

  with  $(x_i)_{i = 1,\dots,n}$ is a subdivision of $\left[a,b\right]$ and $(y_j)_{j = 0,\dots,m}$ a subdivision of 
  $\left[c,d\right]$ we note : $A(\Delta_{ij}) = (x_i - x_{i-1})(y_j -y_{j-1})$

  and $\delta(\sigma) = \max_{i,j} \text{diam}(\Delta_{ij})$

  2) 
  \[\forall i,j \quad M_{ij} = \sup_{(x,y)\in\Delta_{ij}}f(x,y) \quad \text{and} \quad m_{ij} = \inf_{(x,y)\in\Delta_{ij}}f(x,y)\]

  the upper Darboux sum of $f$ with respect to $\sigma$ is :
  \[S(f,\sigma) = \sum_{i=1}^n\sum_{j=1}^m M_{ij}A(\Delta_{ij})\]

  the lower Darboux sum of $f$ with respect to $\sigma$ is :
  \[s(f,\sigma) = \sum_{i=1}^n\sum_{j=1}^m m_{ij}A(\Delta_{ij})\]

  3) if we give $\zeta = \{(\zeta_i^1,\zeta_j^2) \mid \zeta_i^1 \in \left[x_{i-1},x_i\right] \text{ and }\zeta_j^2 \in \left[y_{j-1},y_j\right]\}$

  The Riemann sum with respect to $\sigma$ and $\zeta$ is :
  \[R(f,\sigma,\zeta)= \sum_{i=1}^n\sum_{j=1}^m f(\zeta_i^1,\zeta_j^2)A(\Delta_{ij})\]
  \[R(f,\sigma,\zeta)= \sum_{i=1}^n\sum_{j=1}^m f(\zeta_i^1,\zeta_j^2)(x_i - x_{i-1})(y_j - y_{j-1})\]

\end{definition}

\subsection{Example:}

\[ f(x,y) = \alpha \quad; \forall(x,y) \in \Delta\]
\[M_{ij} = m_{ij} = f(\zeta_i^1, \zeta_j^2) = \alpha\]
\[S(f,\sigma) = s(f, \sigma) = R(f,\sigma, \zeta) = \alpha(b-a)(d-c) = \alpha A(\Delta)\]

\begin{definition}
  Riemann integrability:

  With the same notations above, we put:
  \[S^+(f)= \inf(S(f,\sigma)) \quad \text{and} \quad S^-(f)= \sup (s(f,\sigma))\]

We say that $f$ is Riemann integrable on $\Delta$ if $S^+(f) = S^-(f)$
\[\Leftrightarrow \forall \epsilon > 0; \exists \sigma \quad S(f, \sigma) - s(f, \sigma) < \epsilon\]
and we denote the common value by $\iint_{\Delta}f(x,y)dxdy$
\[\Leftrightarrow \exists(\sigma_n) \text{ of subdivisions of } \Delta / \lim S (f, \sigma_n) = \lim s(f, \sigma_n)\]

\underline{Corollary:} if  $f : \Delta = \left[a,b\right]\times\left[c,d\right]$ is Riemann integrable then:

\[\lim\limits_{n \to +\infty} \frac{(b-a)(d-c)}{n^2}\sum_{i=1}^n\sum_{j =1}^n f(a + \frac{i}{n}(b-a),c + \frac{j}{n}(d-c))\]
\[ = \iint_{\Delta}f(x,y)dxdy\]

\end{definition}

\subsection{Example:}
\[f(x,y) = x^2 + y^2 \quad,\Delta = \left[0,1\right]^2= \left[0,1\right]\times\left[0,1\right]\]


\begin{definition}
  Riemann integrability on a bounded domain:

  Let $D$ be a bounded subset of $\R^2$ and $ f: D \to \R $ a bounded function.

  We say that $f$ is Riemann integrable on $D$ if for any rectangle $\Delta \supset D$, the function :
  \[
    \tilde{f}(x,y) = 
    \begin{cases}
      f(x,y) & \text{if } (x,y) \in D \\
      0 & \text{if } (x,y) \in \Delta \setminus D
    \end{cases}
  \]

  is Riemann integrable, in this case :
  \[ \iint_D f(x,y)dxdy := \iint_{\Delta} \tilde{f}(x,y) dxdy\]

\textbf{Remark:}
Any continuous function is Riemann integrable.

\subsection{Properties:}

1)\underline{Linearity:}
If $f,g: D \to \R$ are Riemann integrable then:
$\forall \alpha,\beta \in \R \quad, \alpha f + \beta g$ is Riemann integrable and 
\[\iint_D(\alpha f + \beta g)(x,y)dxdy = \alpha\iint_Df(x,y)dxdy + \beta\iint_Dg(x,y)dxdy \]

2)

If $f \ge 0$ on $D \implies \iint_D f(x,y) dxdy \ge 0$
which is the volume of $\omega = \{(x,y,z) \in \R^3 : (x,y)\in D \text{ and } 0 \le z \le f(x,y)\}$

So; if $f \ge g \implies \iint_D\left[f(x,y) - g(x,y)\right]dxdy$

$= \text{vol}\{(x,y,z) \in \R^3: (x,y) \in D \text{ and } g(x,y) \le z \le f(x,y)\}$ 

3)\underline{Chasles's relation:}

If $D = D_1 \cup D_2$ such that $D_1 \cap D_2$ is at most a curve then:
\[\iint_Df(x,y)dxdy = \iint_{D_1}f(x,y)dxdy + \iint_{D_2}f(x,y)dxdy \]
\end{definition}

\begin{definition}
  $\text{Vol}(D) = A(D) = \iint_Ddxdy$ \quad and the average of $f$ on D is:
  $\dfrac{1}{\text{Vol}(D)}\iint_Df(x,y)dxdy = \dfrac{\iint_Df(x,y)dxdy }{\iint_D dxdy }$
\end{definition}

\textbf{\underline{Remark:}}

If  $f \equiv C \quad$, the average of $f$ is $C$

\subsection{Example:}

$D$: the triangle with vertices $(0,0),(1,0),(0,1)$
(graph) 
and $f(x,y) = xy $
\[
  \tilde{f}(x,y) = 
  \begin{cases}
    xy  & \text{if } x+y \le 1 \\
    0 & \text{if } x+y > 1
  \end{cases}
\]
$\forall x,y \in \left[0,1\right] \times \left[0,1\right]$

\[\iint_D xy dxdy = \lim\limits_{n \to +\infty} \frac{(1-0)(1-0)}{n^2} \sum_{i=1}^n\sum_{j= 1}^n\tilde{f}(\frac{i}{n},\frac{j}{n})\]

\[\lim\limits_{n \to +\infty}\sum_{i =1}^n \left[\sum_{j =1}^n \frac{i}{n}\times\frac{j}{n}\right]\]
\[U_n = \frac{1}{n^2}\left[\sum_{i=1}^n\sum_{j =1}^n \frac{i}{n}\times\frac{j}{n}\right]\]
\[U_n = \frac{1}{n^4} \sum_{i = 1}^n i\left[\sum_{j=1}^{n-i}j\right] = \frac{1}{n^4}\sum_{i =1}^n i\left(\frac{(n-i)(n-i+1)}{2}\right) \]
\[U_n =  \frac{1}{2n^4} \sum_{i=1}^n i (n^2 + i^2 - 2in + n-i) \]
\[U_n =  \frac{1}{2n^4} \sum_{i = 1}^n \left[i(n^2+n) - i^2(2n+1) +i^3\right]\]
\[U_n = \frac{1}{2n^4} \left[(n^2 + n)\frac{(n+1)n}{2} -(2n + 1)\frac{n(n+1)(2n +1)}{6} + \frac{n^2(n+1)^2}{4}\right] \]
\[U_n  = \frac{1}{4}(1 + \frac{1}{n})^2 - \frac{1}{12}(1 + \frac{1}{n})(2 + \frac{1}{n}) + \frac{1}{8}(1 + \frac{1}{n})^2\]

\[\iint_D xy dxdy = \lim\limits_{n \to + \infty} U_n = \frac{1}{4} - \frac{1}{12}\cdot 2  + \frac{1}{8} = \frac{1}{24}\]

\begin{theorem}
  Fubini's theorem:

  Let $f: D \to \R$ Riemann integrable

  1) if 
  \[D  = \{(x,y) \in \R^2 : a \le x\le b \text{ and } \exists g_1,g_2: \left[a,b\right] \to \R \text{ continuous}
    \mid 
    g_1(x) \le y \le g_2(x), \forall x \in \left[a,b\right]\}
  \]

  then: $\iint_D f(x,y)dxdy = \int_a^b \left[\int_{g_1(x)}^{g_2(x)}f(x,y)dy\right]dx$

  2) If
  \[
    D = \{ (x,y) \in \R^2 : c \le y \le d \text{ and } \exists h_1,h_2: \left[c,d\right] \to \R \text{ continuous}
    \mid h_1(y) \le x \le h_2(y), \forall y \in \left[c,d\right] \}
  \]

  then: $\iint_D f(x,y)dxdy = \int_c^d \left[\int_{h_1(y)}^{h_2(y)}f(x,y)dx\right]dy$
\end{theorem}

\subsection{Remark:}
We can use the Chasles's formula to apply this theorem.

\subsection{Example 01:}
\textbf{Problem:} Show that a disk of radius $R$ centered at the origin has area $\pi R^2$.
 
\textbf{Setup Using Cartesian Coordinates (Quick Version):}
 
The disk is:
\[
D = \{(x,y) \in \mathbb{R}^2 : x^2 + y^2 \leq R^2\}
\]
 
By symmetry (four equal quadrants), we compute:
\[
\text{Area} = 4 \iint_{D_1} dxdy
\]
 
where $D_1$ is the first quadrant portion:
\[
D_1 = \{(x,y) : 0 \leq x \leq R, \, 0 \leq y \leq \sqrt{R^2-x^2}\}
\]
 
\textbf{Compute in Cartesian Coordinates:}
 
\begin{align*}
\text{Area} &= 4\int_0^R \int_0^{\sqrt{R^2-x^2}} dy \, dx \\[0.5em]
&= 4\int_0^R \sqrt{R^2-x^2} \, dx
\end{align*}
 
Using the antiderivative $\int \sqrt{R^2-x^2} \, dx = \frac{x\sqrt{R^2-x^2}}{2} + \frac{R^2}{2}\arcsin\left(\frac{x}{R}\right)$:
 
\begin{align*}
&= 4\left[\frac{x\sqrt{R^2-x^2}}{2} + \frac{R^2}{2}\arcsin\left(\frac{x}{R}\right)\right]_0^R \\[0.5em]
&= 4\left(0 + \frac{R^2}{2} \cdot \frac{\pi}{2} - 0 - 0\right) \\[0.5em]
&= 4 \cdot \frac{\pi R^2}{4} \\[0.5em]
&= \pi R^2
\end{align*}
 
\textbf{Simpler Approach: Using Polar Coordinates:}
 
In polar coordinates $(r, \theta)$ where $x = r\cos\theta$ and $y = r\sin\theta$, the disk becomes:
\[
D = \{(r,\theta) : 0 \leq r \leq R, \, 0 \leq \theta \leq 2\pi\}
\]
 
The area element transforms: $dxdy = r \, dr \, d\theta$
 
Therefore:
\begin{align*}
\text{Area} &= \iint_D dxdy \\[0.5em]
&= \int_0^{2\pi} \int_0^R r \, dr \, d\theta \\[0.5em]
&= \int_0^{2\pi} \left[\frac{r^2}{2}\right]_0^R d\theta \\[0.5em]
&= \int_0^{2\pi} \frac{R^2}{2} \, d\theta \\[0.5em]
&= \frac{R^2}{2} \cdot 2\pi \\[0.5em]
&= \pi R^2
\end{align*}
 
\boxed{\text{Area of disk} = \pi R^2}
 
\textbf{Intuition (Polar Method):} In polar coordinates, we're summing up infinitesimal rings. Each ring at radius $r$ has circumference $2\pi r$ and thickness $dr$, so its area is $2\pi r \, dr$. Integrating from $r=0$ to $r=R$ gives the total area. The factor $r$ in the Jacobian $r \, dr \, d\theta$ accounts for the fact that rings farther from the origin have larger circumference.
 
\subsection{Example 02:}
(trapezoid area proof)
\textbf{Problem:} Prove that the area of a trapezoid with parallel sides of length $A$ and $B$ and height $h$ is $\frac{(A+B)h}{2}$.
 
\textbf{Setup:} Consider a trapezoid in the $xy$-plane where:
\begin{itemize}
  \item The bottom base (at $y=0$) has length $A$
  \item The top base (at $y=h$) has length $B$
  \item The height is $h$
  \item The left side has slope $\frac{a}{h}$ (starting at $x = 0$ when $y = 0$)
  \item The right side has slope $\frac{a+B-A}{h}$ (starting at $x = A$ when $y = 0$)
\end{itemize}
 
\textbf{Domain Description:}
 
For a fixed height $y$ (where $0 \leq y \leq h$), the trapezoid spans horizontally:
\[
D = \left\{ (x,y) \in \mathbb{R}^2 \mid 0 \leq y \leq h, \quad \frac{a}{h}y \leq x \leq \frac{a+B-A}{h}y + A \right\}
\]
 
The left boundary is: $x_{\text{left}}(y) = \frac{a}{h}y$
 
The right boundary is: $x_{\text{right}}(y) = \frac{a+B-A}{h}y + A$
 
\textbf{The Width at Height $y$:}
 
The width of the trapezoid at height $y$ is the distance between left and right boundaries:
\begin{align*}
\text{Width}(y) &= x_{\text{right}}(y) - x_{\text{left}}(y) \\
&= \left(\frac{a+B-A}{h}y + A\right) - \frac{a}{h}y \\
&= \frac{B-A}{h}y + A
\end{align*}
 
\textbf{Notice:} This is linear in $y$!
\begin{itemize}
  \item At $y = 0$: Width $= A$ (the bottom base)
  \item At $y = h$: Width $= \frac{B-A}{h} \cdot h + A = B$ (the top base)
\end{itemize}
 
\textbf{Compute the Area:}
 
\begin{align*}
A(D) &= \int_0^h \left[\int_{\frac{a}{h}y}^{\frac{a+B-A}{h}y + A} dx\right] dy \\[0.5em]
&= \int_0^h \left[\frac{B-A}{h}y + A\right] dy \\[0.5em]
&= \left[Ay + \frac{B-A}{2h}y^2\right]_0^h \\[0.5em]
&= Ah + \frac{B-A}{2h} \cdot h^2 \\[0.5em]
&= Ah + \frac{(B-A)h}{2} \\[0.5em]
&= \frac{2Ah + (B-A)h}{2} \\[0.5em]
&= \frac{(A+B)h}{2}
\end{align*}
 
\boxed{\text{Area of trapezoid} = \frac{(A+B)h}{2}}
 
\textbf{Intuition:} The width of the trapezoid changes linearly from $A$ at the bottom to $B$ at the top. The average width is $\frac{A+B}{2}$, and when multiplied by the height $h$, we get the total area.
 



\subsection{Example 03:}

\textbf{Problem:} Compute $\displaystyle \iint_D xy \, dxdy$ where $D$ is the triangular region in the first quadrant bounded by $x + y = 1$.
 
\textbf{Visualizing the Domain:}
 
The region $D$ is defined by:
\[
D = \{(x,y) \in \mathbb{R}^2 : x \geq 0, \, y \geq 0, \, x + y \leq 1\}
\]
 
This is a right triangle with vertices at $(0,0)$, $(1,0)$, and $(0,1)$.
 
\textbf{Express the Domain Explicitly:}
 
To set up the integral, we choose $x$ as the outer variable:
\begin{itemize}
  \item $x$ ranges from $0$ to $1$
  \item For each fixed $x$, $y$ ranges from $0$ to $1-x$
\end{itemize}
 
Therefore:
\[
D = \{(x,y) \in \mathbb{R}^2 : 0 \leq x \leq 1 \text{ and } 0 \leq y \leq 1-x\}
\]
 
\textbf{Set Up the Double Integral:}
 
\begin{align*}
\iint_D xy \, dxdy &= \int_0^1 x \left[\int_0^{1-x} y \, dy\right] dx
\end{align*}
 
\textbf{Evaluate the Inner Integral:}
 
The inner integral integrates $y$ with respect to $y$:
\begin{align*}
\int_0^{1-x} y \, dy &= \left[\frac{y^2}{2}\right]_0^{1-x} \\
&= \frac{(1-x)^2}{2} - 0 \\
&= \frac{(1-x)^2}{2}
\end{align*}
 
So our integral becomes:
\begin{align*}
\iint_D xy \, dxdy &= \int_0^1 x \cdot \frac{(1-x)^2}{2} \, dx \\
&= \frac{1}{2}\int_0^1 x(1-x)^2 \, dx
\end{align*}
 
\textbf{Expand $(1-x)^2$:}
 
\begin{align*}
(1-x)^2 &= 1 - 2x + x^2
\end{align*}
 
Therefore:
\begin{align*}
x(1-x)^2 &= x(1 - 2x + x^2) \\
&= x - 2x^2 + x^3
\end{align*}
 
\textbf{Evaluate the Outer Integral:}
 
\begin{align*}
\frac{1}{2}\int_0^1 (x - 2x^2 + x^3) \, dx &= \frac{1}{2}\left[\frac{x^2}{2} - \frac{2x^3}{3} + \frac{x^4}{4}\right]_0^1 \\[0.5em]
&= \frac{1}{2}\left(\frac{1}{2} - \frac{2}{3} + \frac{1}{4}\right) \\[0.5em]
&= \frac{1}{2}\left(\frac{6 - 8 + 3}{12}\right) \\[0.5em]
&= \frac{1}{2} \cdot \frac{1}{12} \\[0.5em]
&= \frac{1}{24}
\end{align*}
 
\boxed{\iint_D xy \, dxdy = \frac{1}{24}}
 
\textbf{Intuition:} The function $xy$ is small near the axes (where $x$ or $y$ is near $0$) and largest near the point $(1/3, 1/3)$ inside the triangle. The integral captures this weighted sum.
 


\subsection{Example 04:}
Let us compute volume of $\Omega$ bounded above by the paraboloid $z = 1 - x^2 -y^2$ and below by the plane $z = 1-y$

Compute the volume of the solid $\Omega$ bounded above by the paraboloid $z = 1 - x^2 - y^2$ and below by the plane $z = 1 - y$.
 
 
\textbf{Step 1: Find the Projection (Domain $D$)}
 
The solid is bounded by two surfaces:
\begin{align*}
z &= 1 - x^2 - y^2 \quad \text{(paraboloid)} \\
z &= 1 - y \quad \text{(plane)}
\end{align*}
 
Setting them equal to find where they intersect:
\begin{align*}
1 - x^2 - y^2 &= 1 - y \\
x^2 + y^2 - y &= 0 \\
x^2 + \left(y - \frac{1}{2}\right)^2 &= \frac{1}{4}
\end{align*}
 
This is a circle centered at $\left(0, \frac{1}{2}\right)$ with radius $\frac{1}{2}$.
 
The domain $D$ is:
\[
D = \left\{ (x,y) \in \mathbb{R}^2 \mid x^2 + y^2 - y \leq 0 \right\}
\]
 
This describes the region inside the circle where the plane is below the paraboloid.
 
\textbf{Step 2: Set Up the Solid Region}
 
The solid $\Omega$ is described as:
\[
\Omega = \left\{ (x,y,z) \in \mathbb{R}^3 \mid (x,y) \in D \text{ and } 1-y \leq z \leq 1-x^2-y^2 \right\}
\]
 
\textbf{Step 3: Compute the Volume}
 
The volume is:
\begin{align*}
V &= \iint_D \left[(1-x^2-y^2)-(1-y)\right] \, dxdy \\
&= \iint_D (y - x^2 - y^2) \, dxdy
\end{align*}
 
\textbf{Step 4: Express the Domain Explicitly}
 
From $x^2 + y^2 - y \leq 0$, we can write $x^2 \leq y - y^2 = y(1-y)$.
 
For the region to exist, we need $y(1-y) \geq 0$, which means $0 \leq y \leq 1$.
 
For each fixed $y$, we have:
\[
-\sqrt{y - y^2} \leq x \leq \sqrt{y - y^2}
\]
 
Therefore:
\[
D = \left\{ (x,y) \in \mathbb{R}^2 \mid 0 \leq y \leq 1, \quad -\sqrt{y - y^2} \leq x \leq \sqrt{y - y^2} \right\}
\]
 
\textbf{Step 5: Evaluate the Inner Integral}
 
\begin{align*}
V &= \int_0^1 \left[\int_{-\sqrt{y-y^2}}^{\sqrt{y-y^2}} (y - y^2 - x^2) \, dx\right] dy \\
&= \int_0^1 \left[(y-y^2)x - \frac{1}{3}x^3\right]_{-\sqrt{y-y^2}}^{\sqrt{y-y^2}} dy
\end{align*}
 
By symmetry (the integrand is even in $x$):
\begin{align*}
&= \int_0^1 2\left[(y-y^2)\sqrt{y-y^2} - \frac{1}{3}\left(\sqrt{y-y^2}\right)^3\right] dy \\
&= \int_0^1 2\sqrt{y-y^2}\left[(y-y^2) - \frac{1}{3}(y-y^2)\right] dy \\
&= \int_0^1 2\sqrt{y-y^2} \cdot \frac{2}{3}(y-y^2) \, dy \\
&= \frac{4}{3}\int_0^1 (y-y^2)^{3/2} \, dy
\end{align*}
 
\textbf{Step 6: Trigonometric Substitution}
 
\textbf{Why substitute?} We have $(y - y^2)^{3/2}$. To handle this, we complete the square:
\begin{align*}
y - y^2 &= \frac{1}{4} - \left(y - \frac{1}{2}\right)^2
\end{align*}
 
This suggests a sine substitution. Let:
\begin{align*}
y - \frac{1}{2} &= \frac{1}{2}\sin(t)
\end{align*}
 
Then:
\begin{align*}
y - y^2 &= \frac{1}{4} - \frac{1}{4}\sin^2(t) = \frac{1}{4}(1 - \sin^2(t)) = \frac{1}{4}\cos^2(t) \\
\sqrt{y-y^2} &= \frac{1}{2}|\cos(t)| = \frac{1}{2}\cos(t) \quad \text{(for } t \in [-\pi/2, \pi/2])\\
dy &= \frac{1}{2}\cos(t) \, dt
\end{align*}
 
\textbf{Bounds:} 
- When $y = 0$: $\sin(t) = -1 \Rightarrow t = -\pi/2$
- When $y = 1$: $\sin(t) = 1 \Rightarrow t = \pi/2$
 
\textbf{Step 7: Transform the Integral}
 
\begin{align*}
V &= \frac{4}{3}\int_{-\pi/2}^{\pi/2} \left(\frac{1}{2}\cos(t)\right)^3 \cdot \frac{1}{2}\cos(t) \, dt \\
&= \frac{4}{3}\int_{-\pi/2}^{\pi/2} \frac{1}{16}\cos^4(t) \, dt \\
&= \frac{1}{12}\int_{-\pi/2}^{\pi/2} \cos^4(t) \, dt
\end{align*}
 
\textbf{Step 8: Evaluate $\int_{-\pi/2}^{\pi/2} \cos^4(t) \, dt$}
 
Using the power-reduction formula for even powers:
\[
\cos^4(t) = \left(\cos^2(t)\right)^2 = \left(\frac{1 + \cos(2t)}{2}\right)^2 = \frac{1 + 2\cos(2t) + \cos^2(2t)}{4}
\]
 
For $\cos^2(2t)$:
\[
\cos^2(2t) = \frac{1 + \cos(4t)}{2}
\]
 
Therefore:
\begin{align*}
\cos^4(t) &= \frac{1}{4}\left(1 + 2\cos(2t) + \frac{1 + \cos(4t)}{2}\right) \\
&= \frac{1}{4}\left(\frac{3}{2} + 2\cos(2t) + \frac{1}{2}\cos(4t)\right) \\
&= \frac{3}{8} + \frac{1}{2}\cos(2t) + \frac{1}{8}\cos(4t)
\end{align*}
 
Now integrate:
\begin{align*}
\int_{-\pi/2}^{\pi/2} \cos^4(t) \, dt &= \int_{-\pi/2}^{\pi/2} \left(\frac{3}{8} + \frac{1}{2}\cos(2t) + \frac{1}{8}\cos(4t)\right) dt \\
&= \left[\frac{3t}{8} + \frac{1}{4}\sin(2t) + \frac{1}{32}\sin(4t)\right]_{-\pi/2}^{\pi/2} \\
&= \frac{3\pi}{16} + 0 + 0 - \left(-\frac{3\pi}{16} + 0 + 0\right) \\
&= \frac{3\pi}{8}
\end{align*}
 
\textbf{Step 9: Final Answer}
 
\begin{align*}
V &= \frac{1}{12} \cdot \frac{3\pi}{8} \\
&= \frac{3\pi}{96} \\
&= \frac{\pi}{32}
\end{align*}
 
\boxed{V = \frac{\pi}{32}}

\end{document}


\newpage
\subsection{Exercises Series 1}
\vspace*{0.2cm}
\Large{\underline{\textbf{Exercise 1}}}
\normalsize{}\\
\noindent Using Riemann sums, calculate $\displaystyle \iint_D(x + 2y)\ dxdy,$ where $D = \{(x, y)\in \R^2:\ x\geq 0,\ y\geq 0,\ x + y = 1 \}$

\vspace*{0.34cm}
\noindent\Large{\underline{\textbf{Exercise 2}}}
\normalsize{}
\begin{enumerate}
\item Calculate the following integrals:
        \begin{align*}
            &I_1 = \iint_D \dfrac{y}{\sqrt{1 + xy + y^2}}\ dxdy, 
            \qquad D = \{(x, y)\in \R^2:\ -2\leq x \leq 2,\ 3\leq y \leq 7\}\text{.}\\
            &I_2 = \iint_D (x + 3y)\ dxdy, \qquad \qquad \ \ D = \{(x, y)\in \R^2:\ x \geq 0,\ x^2 + y^2 \leq 1\}\text{.}
        \end{align*}
\item Let $I$ be an integral defined as $$\displaystyle \int_0^1 \dfrac{\ln(1 + x)}{1 + x^2}\ dx.$$
        $i)$ Show that for any real number $x \geq 1,$ we have $$\ln (1 + x) = \int_0^1 \dfrac{x}{1 + xy}\ dy.$$\\
        $ii)$ Prove that $\displaystyle I = \iint_{[0, 1]^2} \dfrac{x}{(1 + xy)(1 + x^2)}\ dxdy$, and that
        $\displaystyle 2I = \iint_{[0, 1]^2} \dfrac{x + y}{(1 + x^2)(1 + y^2)}\ dxdy.$\\

        $iii)$ Deduce the value of $I.$
\end{enumerate}

\vspace*{0.34cm}
\noindent\Large{\underline{\textbf{Exercise 3}}}
\normalsize{}\\
\noindent $i)$ Find the volume of the first octant (i.e. $(\R^+)^3$) bounded by the graphs of the equations:
\begin{align*}
    \textbf{1.}& \qquad x^2 + z^2 = 1,\quad x^2 + y^2 = 1.\\
    \textbf{2.}& \qquad x + y = z,\quad x^2 + y^2 = 4.
\end{align*}

\noindent $ii)$ Do the same for the solid $\Omega$ in each of the following cases:
\begin{align*}
    \textbf{1.}& \qquad x = 0,\quad y = 0,\quad x + y = 1,\quad x^2 + x + 1 = z.\\
    \textbf{2.}& \qquad 16 - 4x^2 - y^2 = 4z,\quad x^2 + y^2 = 2x.
\end{align*}

\noindent\Large{\underline{\textbf{Exercise 4}}}
\normalsize{}\\
\noindent Given $\rho$ as the density function, find the center of mass, $(x_c, y_c, z_c)$, and the moments of inertia:
    $$\Omega = \left\{(x,y,z)\in\R^3:\ x^2 + y^2 \leq z \leq \sqrt{4 - x^2 - y^2}\right\},\ \qquad \qquad \rho(x,y,z) = kz.$$

\noindent\Large{\underline{\textbf{Exercise 5}}}
\normalsize{}\\
\noindent Find the area of the surface given by the equation $z = f(x, y)$ over $D$ in the following cases.
\begin{align*}
    \textbf{1.}& \qquad f\ :\ (x, y) \longmapsto 8 + 2x - 3y, \qquad \qquad D = \{(x, y)\in \R^2:\ x^2 + y^2 \leq 9\}\\
    \textbf{2.}& \qquad f\ :\ (x, y) \longmapsto -\ln(\sin(x)), \qquad \qquad D = \left\{(x, y)\in \R^2:\ 
                                                                  0 \leq x \leq \frac{\pi}{4},\ 0 \leq y \leq \tan (x)\right\}
\end{align*}

\vspace*{0.34cm}
\noindent\Large{\underline{\textbf{Exercise 6}}}
\normalsize{}\\
\noindent Use a triple integral to find the volume of the solid bounded by the graphs of the equations below:
\begin{align*}
    \textbf{1/ }& \qquad z = 9 - x^3,\quad y = 2 - x^2,\quad y = 0,\quad z = 0,\quad x \geq 0,\\
    \textbf{1/ }& \qquad z = 2 - y,\quad z = 4 - y^2,\quad x = 0,\quad x = 3,\quad y = 0,\\
\end{align*}

\pagebreak
\vspace*{0.34cm}
\noindent\Large{\underline{\textbf{Exercise 7}}}
\normalsize{}\\
\noindent Use a change of variables to find the volume of the region lying below the surface $z = f (x, y)$ and above the plane region $D$.
\begin{align*}
    \textbf{1/}& \ f\ :\ (x, y) \longmapsto (3x + 2y)^2 \sqrt{2y -x}, \qquad \text{$D$ is the parallelogram with vertices
                                                                $(0, 0),\ (-2, 3),\ (2, 5),$ and $(4, 2)$.}\\
    \textbf{2/}& \ f\ :\ (x, y) \longmapsto \dfrac{xy}{1 + x^2y^2}, \qquad \qquad \qquad \ \text{$D$ is bounded by the graphs of $xy = 1,\ xy = 4,\ x = 1,$ and $x = 4$.}
\end{align*}
\vspace*{0.34cm}

\noindent\Large{\underline{\textbf{Exercise 8}}}
\normalsize{}\\
\noindent Find the mass, the center of mass, and the moments of inertia $I_x$ and $I_y$ of
                                the lamina bounded by the graphs of the equations with the given density $\rho$ :
\begin{align*}
    \textbf{1. }& \qquad y = \dfrac{4}{x},\quad y = 0,\quad x = 1,\quad x = 4, \qquad \rho(x,y) = kx^2.\\
    \textbf{2. }& \qquad y = 4 - x^2,\quad y = 0, \qquad \qquad \qquad\ \rho(x,y) = ky.\\
    \textbf{3. }& \qquad y = \sqrt{a^2 - x^2},\quad 0 \leq y \leq x, \qquad \quad \rho(x,y) = k.
\end{align*}
\noindent\Large{\underline{\textbf{Exercise 9}}}
\normalsize{}\\
\noindent $i)$ Integrate the given function over the indicate region $D$.
\begin{align*}
    \textbf{1. }& f(x,y,z) = 8xyz,     \qquad \qquad \text{ $D$ is bounded by $y = x^2,\ y + z = 9,\ z = 0$.} \\
    \textbf{2. }& f(x,y,z) = yz + x^2, \qquad \quad \text{$D$ is bounded
                                                    by $x= 0,\ y = 0,\ z = 0,$ and $x + y + 2z \leq 1$.} \\
    \textbf{3. }& f(x,y,z) = 1 - z^2,  \qquad \quad \ \text{ $D$ is
                                the tetrahedron with vertices $(0,0,0),\ (1,0,0),\ (0,2,0)$, and $(0,0,3)$.}
\end{align*}
\noindent $ii)$ Using an appropriate change of variables, calculate the triple integral of the 3-valued functions below:
\begin{align*}
    \textbf{1. }& f(x,y,z) = x^2 + y^2,  \qquad \qquad \quad \text{$D = 
                                                        \{(x,y,z)\in \R^3:\ x^2 + y^2 \leq (z - 1)^2,\ 0\leq z\leq 1\}$.}\\
    \textbf{2. }& f(x,y,z) = z\cos(x^2 + y^2),  \qquad \text{$D = 
                                                        \{(x,y,z)\in \R^3:\ x^2 + y^2 + z^2 \leq 1,\ z \leq 0\}$.}
\end{align*}

\vspace*{0.34cm}
\noindent\Large{\underline{\textbf{Exercise 10}}}
\normalsize{}\\
\noindent $i)$ Following the indicated change of variables, calculate $\displaystyle \iint_{D_1} xy\ dxdy,$ where
$$D_1 = \{(x,y)\in\R^2:\ x\geq 0,\ y\geq 0,\ x^{2/3} + y^{2/3} \leq 1\},\qquad \text{  and} \qquad
                                                       x = r\cos^3(\theta),\quad y = r\sin^3(\theta).$$
\noindent $2)$ Calculate, using the same approach, $\displaystyle \iint_{D_2} (x^2 + y^2)(x^2 - y^2)\ dxdy,$ where
$$D_2 = \{(x,y)\in\R^2:\ 0 < x < y,\ 0 < a < xy < b, y^2 - x^2 < 1\},\qquad \text{  and} \qquad 
                                                       u = xy,\quad v = y^2 - x^2.$$\\
\noindent\Large{\underline{\textbf{Exercise 11}}}
\normalsize{}\\
\noindent $i)$ Integrate the given function $f$ over the indicated region $D$.
\begin{align*}
    \textbf{1. }& \ f(x,y,z) = 2x - y + z, \qquad \text{ $D$ is bounded by $z = x^2 + y^2,\ z = 0,\ 
                                                                    x = 0,\ x = 1,\ y = -2$, and $y = 2$.}\\
    \textbf{2. }& \ f(x,y,z) = \dfrac{z}{\sqrt{x^2 + y^2}},\ \qquad D = \{(x,y,z)\in\R^3:\ x^2 + y^2 + z^2 \leq 1,\
                                                              y^2 - 2xz \leq 0,\ 4z^2 \leq x^2 + y^2,\ z\geq 0\}
\end{align*}
\noindent $ii)$ Let $\alpha > 0$, and $\Omega = \{(x,y,z)\in\R^3:\ x^2 + y^2 + z^2 \leq 1\}$. Calculate
$$\displaystyle \iiint_\Omega \frac{1}{\sqrt{x^2 + y^2 + (z - \alpha)^2}}\ dxdydz$$


\newpage
\section{Vectorial Analysis}
% \input{Chapters/Chapter 2.tex}
\newpage
\subsection{Exerices Series 2}
\vspace*{0.2cm}
\Large{\underline{\textbf{Exercise 1}}}
\normalsize{}
\begin{enumerate}
    \item Find the curl of $F$ in reach of the following cases:
    \begin{align*}
        i)\quad  F\ :\ (x,y,z) &\longmapsto 
                        \dfrac{yz}{y - z}\vec{i} + \dfrac{xz}{x - z}\vec{j} + \dfrac{xy}{x - y}\vec{k}\\
        ii)\quad  F\ :\ (x,y,z) &\longmapsto 
                        \sin(x - y)\vec{i} + \sin(y - z)\vec{j} + \sin(z - x)\vec{k}
    \end{align*}
    \item Determine whether the vector field is conservative and if so, find a potential function for it.
    \begin{align*}
        i)\quad  F\ :\ (x,y,z) &\longmapsto 
                                    xy^2z^2\ \vec{i} + x^2yz^2\ \vec{j} + x^2y^2z\ \vec{k}\\
        ii)\quad  F\ :\ (x,y,z) &\longmapsto 
                                    \sin(z)\vec{i} + \sin(x)\vec{j} + \sin(y)\vec{k}
    \end{align*}
    \item In the following, find the divergence of the vector field $F$:
    \begin{align*}
        &i)\quad  F\ :\ (x,y) \longmapsto xe^x\ \vec{i} + ye^y\ \vec{j}\\
        i&i)\quad  F\ :\ (x,y,z) \longmapsto \ln(x^2 + y^2)\vec{i} + xy\ \vec{j} + \ln(y^2 + z^2)\vec{k}
    \end{align*}
\end{enumerate}
\vspace*{0.34cm}
\Large{\underline{\textbf{Exercise 2}}}
\normalsize{}\\
\noindent Let $F$ and $G$ be two vector fields and $f$ a scalar field. Prove the following properties:
    \begin{align*}
        &1.) \quad \nabla\cdot(F\times G) = (\nabla\times F)\cdot G - F\cdot(\nabla\times G) \\
        &2.) \quad \nabla\cdot(f\cdot G) = f\cdot \nabla G + \nabla f \cdot G \\
        &3.) \quad \nabla\cdot(\nabla\times F) = 0, \qquad \nabla\times(\nabla f) = 0 \\
        &4.) \quad \nabla\times(f\cdot G) = f\cdot (\nabla\times G) + G\cdot (\nabla f)
    \end{align*}
where $\cdot$ denotes the scalar product, and $\times$ denotes the vector product.\\ \linebreak
\Large{\underline{\textbf{Exercise 3}}}
\normalsize{}
\begin{enumerate}
    \item Evaluate $\displaystyle \int_C(x^2 + y^2)\ d\gamma$, where:\\
        $i)\ C$ is the line segment from $(0,0)$ to $(1,1)$.\\
        $ii)\ C$ is the line segment from $(0,0)$ to $(2,4)$.\\
        $iii)\ C$ is the counterclockwise line around the circle $x^2 + y^2 = 1$ from $(1,0)$ to $(0,1)$.
    \item Evaluate $\displaystyle \int_C(2x + y^2 - z)\ dr$, where:\\
        $i)\ C$ is the broken line connecting $(0,0,0),\ (1,0,0),\ (1,0,1)$, and $(1,1,1)$.\\
        $ii)\ C$ is triangle with vertices $(0,0,0),\ (0,1,0)$, and $(0,1,1)$.\\
\end{enumerate}

\noindent\Large{\underline{\textbf{Exercise 4}}}
\normalsize{}
\begin{enumerate}
\item Evaluate $\displaystyle \int_C F\cdot dr,$ where:
\begin{align*}
    &i)  \ F\ :\ (x,y) \longmapsto xy\ \vec{i} + y\ \vec{j},\qquad \qquad \qquad 
                            \text{ and }\ \ C = \{(x,y)\in(\R^+)^2:\ x^2 + y^2 = 4\}.\\ 
    &ii) \ F\ :\ (x,y,z) \longmapsto xy\ \vec{i} + xz\ \vec{j} + yz\ \vec{k},\quad\ \text{ and }\ \ C\ :\ t \longmapsto
                            t\ \vec{i} + t^2\ \vec{j} + 2t\ \vec{k},\ t\in[0,1].\\
    &iii)\ F\ :\ (x,y,z) \longmapsto x^2\ \vec{i} + y^2\ \vec{j} + z^2\ \vec{k},\quad\ \text{ and }\ \ C\ :\ t \longmapsto
                            2\sin (t)\vec{i} + 2\cos (t)\vec{j} + \frac{t^2}{2}\ \vec{k},\ t\in[0,\pi].
\end{align*}
\item Evaluate $\displaystyle \int_C(x^2 + y^2)\ dx + 2xy\ dy$, where:
\begin{align*}
    & i)  \ C\ :\ t \longmapsto t^3\ \vec{i} + t^2\ \vec{j},\ t\in[0,2]
    & ii) \ C\ :\ t \longmapsto 2\cot t\ \vec{i} + 2\sin t\ \vec{j},\ t\in\left[0,\pi\right]
\end{align*}
\item Evaluate $\displaystyle \int_C yz\ dx + xz\ dy + xy\ dz$, where:
\begin{align*}
    & i)  \ C\ :\ t \longmapsto t\ \vec{i} + 2\ \vec{j} + t\ \vec{k},\ t\in[0,4]
    & ii) \ C\ :\ t \longmapsto t^2\ \vec{i} + t\ \vec{j} + \ t^2\ \vec{k},\ t\in[0,2]
\end{align*}
\end{enumerate}
\newpage
\noindent\Large{\underline{\textbf{Exercise 5}}}
\normalsize{}\\
\noindent Use Green’s Theorem to evaluate the indicated line integral:
\begin{align*}
    &\textbf{1)}\ \int_C 2xy\ dx + (x + y)\ dy,\ 
              \text{ where $C$ is the region laying between $y = 0$ and $y = 1 - x^2$.}\\
    &\textbf{2)}\ \int_C 2\arctan\left(\dfrac{y}{x}\right)\ dx + \ln(x^2 + y^2)\ dy,\ 
              \text{ where $C$ is the eclipse bounded by $(x - 4)^2 + 4(y - 4)^2 = 4$.}\\
    &\textbf{3)}\ \int_C (x^2 - y^2)\ dx + 2xy\ dy,\ 
              \text{ where $C = \{(x,y)\in\R^2:\ x^2 + y^2 \leq 16\}$.}
\end{align*}

\noindent\Large{\underline{\textbf{Exercise 6}}}
\normalsize{}
\begin{enumerate}
    \item Find the area of the torus $T : (x,y) \longmapsto 
                      (2 + \cos x)\cos y\ \vec{i} + (2 + \cos x)\sin y\ \vec{j} + \sin x\ \vec{k}$.
    \item Find the vector field $F$ whose graph is the indicated surface:
        \begin{align*}
            &i)\   \ z = y.
            &ii)\   \ z = \sqrt{4x^2 + 9y^2}.\\
            &iii)\ \ 4x^2 + y^2 = 16.
            &iiii)\ \ \frac{x^2}{9} + \frac{y^2}{4} + z^2 = 1.
        \end{align*}
\end{enumerate}
\noindent\Large{\underline{\textbf{Exercise 7}}}
\normalsize{}
\begin{enumerate}
    \item Evaluate $\displaystyle\iint_S f(x,y,z)\ d\sigma$, where:
        \begin{align*}
            i)\ f\ &:\ (x,y,z) \longmapsto x^2 + y^2 + z^2,\ \ \
                            \text{and $S = \{(x,y,z)\in\R^3:\ z = x + y,\ x^2 + y^2 \leq 1$\}.}\\
            ii)\ f\ &:\ (x,y,z) \longmapsto \sqrt{x^2 + y^2 + z^2},\
                            \text{and $S = \{(x,y,z)\in\R^3:\ z = \sqrt{x^2 + y^2},\ (x - 1)^2 + y^2 \leq 1$\}.}
        \end{align*}
    \item Find the flux of the vector field through the closed surfaces $S$:
        \begin{align*}
            i)\ f\ :\ (x,y,z) \longmapsto 3z\ \vec{i} - 4\vec{j} + y\ \vec{k},\ 
                                        \text{ and $S = \{(x,y,z)\in(\R^+)^3:\ z = 1 - x - y\}$.}\\
            ii)\ f\ :\ (x,y,z) \longmapsto x\ \vec{i} + y\ \vec{j} + z\ \vec{k},\
                                        \text{ and $S = \{(x,y,z)\in(\R^+)^3:\ z = 1 - x - y\}$.}
        \end{align*}
\end{enumerate}

\noindent\Large{\underline{\textbf{Exercise 8}}}
\normalsize{}\\
\noindent Use the Divergence theorem to compute the flux of $F : (x,y,z) \longmapsto 
x^2\ \vec{i} - 2xy\ \vec{j} + xyz^2\ \vec{k}$, trough the surface $$S = \{(x,y,z)\in\R^3:\ z = \sqrt{a^2 - x^2 - y^2}\}.$$

\noindent\Large{\underline{\textbf{Exercise 9}}}
\normalsize{}\\
Check the Stokes theorem for $F : (x,y,z) \longmapsto x^2\ \vec{i} + y^2\ \vec{j} + z^2\ \vec{k}$
trough the surface\\ $$S = \{(x,y,z)\in\R^3:\ z = y^2,\ 0\leq x\leq a,\ 0\leq y\leq a\}.$$


% \newpage
% \input{Chapters/Chapter 3.tex}
% \subsection{Exercises Series 3}
\vspace*{0.2cm}
\Large{\underline{\textbf{Exercise 1}}}
\normalsize{}
\begin{enumerate}
    \item Using the definition, study the nature and calculate the possible sum 
    of the series $\displaystyle\sum u_n $ in each case:
        $$ u_n = \dfrac{\cos \alpha n}{2^n},\ \alpha\in\R,\qquad \qquad \qquad 
           u_n = \dfrac{n + 1}{3^n}, \qquad \qquad \qquad u_n = \dfrac{n - (-1)^n}{3^n}, $$
        $$ u_n = \dfrac{n}{n^4 + n^2 + 1},\qquad \qquad 
           u_n = \dfrac{1 + 2 + 3 + \cdot + n}{1^3 + 2^3 + 3^3 + \cdot + n^3},\qquad \qquad 
           u_n = \ln(1 - \dfrac{1}{n^2}). $$
    \item Let $\alpha$ and $\beta$ be real numbers, and for any $n\in\N$, we have 
    $u_n = \ln(n) + \alpha\ln(n + 2) + b\ln(n + 3).$ Find the values of $\alpha$ and $\beta$ for which the series
    $\displaystyle\sum u_n$ converges and, in that case, calculate its sum.\\
\end{enumerate}

\Large{\underline{\textbf{Exercise 2}}}
\normalsize{}
\begin{enumerate}
    \item Indicate the nature of the series in each case:
        $$ u_n = \sqrt{n^2 + n} - n, \qquad\qquad u_n = \dfrac{\ln n}{n^2},\qquad \qquad
           u_n = \dfrac{1 + \ln n}{\sqrt{n}},$$
        $$ u_n = \dfrac{n + \sqrt{n}}{2n^3 - 1}, \qquad
           u_n = \left(\sqrt[n]{2} + \sqrt[n]{3}\right)^{-n^2}, \qquad
           u_n = \dfrac{(n + 1)(n + 2)\cdot(2n)}{(2n)^n}$$
    \item Study the nature of the following series:
        $$ \sum_{n\in\N}\dfrac{1}{\mathrm{C}^n_{2n}},\qquad 
           \sum_{n\N^*}\left(1 - \dfrac{n}{2}\ln\left(\dfrac{n + 1}{n - 1}\right)\right),\qquad
           \sum_{n\in\N}\left(\sqrt[3]{n^3 + \alpha n} - \sqrt{n^2 + 3}\right),\ \alpha\in\R $$
\end{enumerate}

\noindent\Large{\underline{\textbf{Exercise 3}}}
\normalsize{}\\
\noindent Study the nature of the following series:
$$ \sum_{n\in\N^*}\dfrac{\cos n}{n^2}, \qquad \qquad \sum_{n\in\N^*}\frac{\cos n}{n},\qquad \qquad 
   \sum_{n\in\N^*}\dfrac{\cos^2n}{n}. $$

\noindent\Large{\underline{\textbf{Exercise 4}}}
\normalsize{}\\
\noindent Let $f : [0, +\infty] \longmapsto \R$ be a continuous, decreasing function, which verifies 
$\displaystyle\lim_{x\longrightarrow +\infty} f(x) = 0.$
\begin{enumerate}
    \item Study using the Leibniz criterion the nature of the series from the general term
            $$ u_n = \int_{\pi n}^{\pi(n + 1)} f(x)\sin x\ dx$$
    \item Deduce the nature of the series of the general term
            $$ u_n = \int_{\sqrt{\pi n}}^{\sqrt{\pi(n + 1)}}\sin (x^2)\ dx$$
\end{enumerate}


% \newpage
% \input{Chapters/Chapter 4.tex}
% \subsection{Exercise Series 3}
\vspace*{0.2cm}
\Large{\underline{\textbf{Exercise 1}}}


% \newpage
% \input{Chapters/Chapter 5.tex}
% \subsection{Exercises Series 5}
\vspace*{0.2cm}
\noindent\Large{\underline{\textbf{Exercise 1}}}
\normalsize{}
\begin{enumerate}
    \item Write an equivalent series with the index of summation beginning at $n = 1$ :
    $$  \sum_{n \in \mathbb{N}} \dfrac{x^{2n+1}}{(2n+1)!}, \qquad \sum_{n \in \mathbb{N}^*} (-1)^n (n+1)\, x^n, \qquad 
        \sum_{n \in \mathbb{N}} \dfrac{(-1)^n}{2n+1} x^{2n+1}.  $$
    \item Determine the radius of convergence of the following power series:
    $$  \sum_{n \in \mathbb{N}} (\sin n)\, x^n, \qquad \sum_{n \in \mathbb{N}} \dfrac{x^n}{(n+1)^{a+1} 3^n}, \qquad 
        \sum_{n \in \mathbb{N}} \arccos\!\left(1 - \dfrac{1}{n^2}\right)x^n, \qquad 
        \sum_{n \in \mathbb{N}} \dfrac{n^n}{n!}\, x^n. $$
    \item Find the interval of convergence of the following power series
    $$  \sum_{n \in \mathbb{N}} \dfrac{(-1)^n}{2^n} (x+1)^n, \qquad 
        \sum_{n \in \mathbb{N}^*} \dfrac{x^n}{n}, \qquad \sum_{n \in \mathbb{N}} \dfrac{(-1)^n\, n!}{3^n} (x-5)^n. $$
\end{enumerate}

\noindent\Large{\underline{\textbf{Exercise 2}}}
\normalsize{}\\
\noindent Determine the radius of convergence as well as the sum of the following power series:
    $$ \sum_{n \in \mathbb{N}} \dfrac{x^n}{(n+1)(n+3)}, \quad 
       \sum_{n \in \mathbb{N}} \left(n^2 - n - 3\right) 3^{n-1} x^n, \quad 
       \sum_{n \in \mathbb{N}} \cosh(na)\, x^n, \quad 
       \sum_{n \in \mathbb{N}} \dfrac{n^2 - n + 4}{n+1}\, x^n, 
       \quad \sum_{n \in \mathbb{N}} \dfrac{x^n}{2n-1}. $$

\noindent\Large{\underline{\textbf{Exercise 3}}}
\normalsize{}\\
\noindent Expand the following functions into power series at the origin (Maclaurin series).
    $$ f_1 : x \longmapsto \dfrac{1}{(3+x)^2}, \qquad 
       f_2 : x \longmapsto x\ln\!\left(x + \sqrt{x^2+1}\right), \qquad 
       f_3 : x \longmapsto \arctan\!\left(\dfrac{1-x^2}{1+x^2}\right). $$

\noindent\Large{\underline{\textbf{Exercise 4}}}
\normalsize{}\\
\noindent Use the definition of Taylor series to find the Taylor series, 
centered at $x_0$ for the function $f$ in the following cases :
    $$  f_1 : x \longmapsto \cos x,\ \ x_0 = \dfrac{\pi}{4}; \qquad 
        f_2 : x \longmapsto e^x,\ \ x_0 = 1; \qquad 
        f_3 : x \longmapsto \sqrt{x},\ \ x_0 = 4.   $$

\noindent\Large{\underline{\textbf{Exercise 5}}}
\normalsize{}\\
\begin{enumerate}
    \item Solve in $\mathbb{R}$ the equation $\displaystyle\sum_{n=0}^{+\infty} (3n+1)^2\, x^n = 0$.
    \item Using integration by parts to establish equality
        $$\iint_{[0,1]^2} xy e^{xy}\, dx\, dy = e - 1 - \sum_{n=1}^{+\infty} \dfrac{1}{n\,(n!)}. $$
\end{enumerate}

\noindent\Large{\underline{\textbf{Exercise 6}}}
\normalsize{}\\
\noindent Let $(a_n)_{n \in \mathbb{N}}$ the sequence defined by $a_0 = 1$ and 
$2a_{n+1} = \displaystyle\sum_{k=0}^{n} \mathrm{C}_n^k\, a_k\, a_{n-k}$, 
we put $f : x \longmapsto \displaystyle\sum_{n=0}^{+\infty} \dfrac{a_n}{n!} x^n$.
\begin{enumerate}
    \item Show that $\forall n\in\N,\ a_n \leq n!$. Deduce that the radius of convergence of 
          $\displaystyle\sum_{n \in \mathbb{N}} \dfrac{a_n}{n!} x^n$ is strictly positive.
    \item Calculate $f'$ as a function of $f$, deduce $f$.
    \item Deduce $a_n$ as function of $n$.
\end{enumerate}


% \newpage
% \input{Chapters/Chapter 6.tex}
% \subsection{Exercises Series 6}
\vspace*{0.2cm}
\noindent\Large{\underline{\textbf{Exercise 1}}}
\normalsize{}


% \input{Chapters/Chapter 7.tex}
% \subsection{Exercises Series 7}
\vspace*{0.2cm}
\noindent\Large{\underline{\textbf{Exercise 1}}}
\normalsize{}\\
\begin{enumerate}
    \item Proof that if $\displaystyle f\in\mathcal{L}^1(\mathbb{R})$ and 
          $\displaystyle\int_{\mathbb{R}} |tf(t)|\, dt$ converges then $\mathcal{F}f$ is derivable. Furthermore,
          $$(\mathcal{F}f)'(x) = -i\mathcal{F}(t \longmapsto tf(t))(x),\ \ \forall x\in\R.$$
    \item Let $\displaystyle\Gamma : t \longmapsto e^{-\frac{t^2}{2}}$. Compute $\mathcal{F}(\Gamma)$.
    \item Let $s > 0$. We consider the function $\displaystyle\Gamma_s : t \longmapsto 
          \dfrac{1}{\sqrt{s}}\, e^{-\frac{t^2}{2s}}.$\\
        \noindent$i)$ Show that $\displaystyle\mathcal{F}(\Gamma_s) = \dfrac{1}{\sqrt{s}}\,\Gamma_{\frac{1}{s}}.$\\
        \noindent$ii)$ Deduce that $\displaystyle\mathcal{F}^2(\Gamma_s) := \mathcal{F}(\mathcal{F}(\Gamma_s))$.\\
        \noindent$iii)$ Let $f\in\mathcal{L}^1(\R,\R)$. Using $\displaystyle\int_{\R}f\mathcal{F}(g)=\int_{\R}g\mathcal{F}(f)$ 
                 and the time shift property,\\ 
                 \indent\quad establish a relationship between $\displaystyle\mathcal{F}^2(f*\Gamma_s)$ and $f*\Gamma_s$.
    \item For any $n\in\mathbb{N}^*$, we define $\gamma_n := \Gamma_{\frac{1}{2n}}$.\\
          $i)$ Show that the sequence $(\gamma_n)_{n \in \mathbb{N}^*}$ is an approximation of unity.\\
             \indent\quad (That is:$\displaystyle\ \forall n\in\N^*,\ \int_{\R} \gamma_n = 1$ and $\forall\varepsilon > 0$, 
             $\displaystyle\ \lim_{n\to\infty} \int_{\mathbb{R} - ]-\varepsilon,\varepsilon[} \gamma_n = 0$).\\
          $ii)$ Using the fact that for any $f\in\mathcal{L}^1(\R, \R)$ the sequence $(f * \gamma_n)_{n \in \N^*}$ 
          converges to $f$ in $\mathcal{L}^1(\R, \R)$, deduce (using the Riesz-Fischer theorem) the Fourier inversion formula
          for a continuous function on $\mathcal{L}^1(\mathbb{R}, \mathbb{R})$ whose Fourier transform is also in 
          $\mathcal{L}^1(\R, \R)$.\\
\end{enumerate}

\noindent\Large{\underline{\textbf{Exercise 2}}}
\normalsize{}\\
\noindent Let $\Pi$ the function defined by
$$\Pi(t) = \begin{cases} 1 & \text{if} \quad 2|t| \leq 1 \\ 0 & \text{if} \quad 2|t| > 1 \end{cases}$$
\begin{enumerate}
    \item Compute the Fourier transform of $\Pi$.
    \item Using the properties of Fourier transform, calculate the Fourier transform of the following functions:
        $$i)\quad t \longmapsto \Pi\!\left(\dfrac{t-1}{2}\right), \qquad ii)\quad t \longmapsto t\ \Pi(t), \qquad 
            iii)\quad t \longmapsto t^2\ \Pi(t).$$
\end{enumerate}

\noindent\Large{\underline{\textbf{Exercise 3}}}
\normalsize{}\\
\noindent We consider the function
$$f(t) = \begin{cases} 1 - |t| & \text{if} \quad |t| \leq 1 \\ 0 & \text{if} \quad |t| > 1 \end{cases}$$
\begin{enumerate}
    \item Proof that $f \in \mathcal{L}^1(\R, \R)$ and compute $\mathcal{F}f$.
    \item Apply the Fourier inversion formula to $f$ at any point in $\R_+$.
    \item Deduce the value of the integral 
            $\displaystyle\int_0^{+\infty} \dfrac{1 - \cos x}{x^2}\cos(\alpha x)\, dx$, $\alpha\in\R_+$.\\
    \item Give the value $\displaystyle\int_0^{+\infty} \dfrac{\sin^2 x}{x^2}\, dx.$
\end{enumerate}
\newpage
\vspace*{0.2cm}
\noindent\Large{\underline{\textbf{Exercise 4}}}
\normalsize{}\\
\noindent For $\alpha > 0$, we put $$f_\alpha : t \longmapsto e^{-\alpha|t|}.$$
\begin{enumerate}
    \item Compute the Fourier transform of $f_\alpha$.
    \item Using the Fourier inversion formula, calculate the value of
          $$\int_0^\infty \dfrac{\cos(xt)}{x^2 + \alpha^2}\ dx, \qquad \text{ for any } t\in\R.$$
    \item Find $a$ and $b$ such that
          $$ \dfrac{1}{\left(x^2 + 1\right)\left(x^2 + 4\right)} = \dfrac{a}{x^2+1} + \dfrac{b}{x^2+4}, \qquad 
          \text{for any } x \in \R.$$
    \item Solve the differential equation
          $$-y'' + y = f_2,$$
          where $y,\, y',\, y''$ and $y$ is derivable on $\R$.\\
\end{enumerate}

\noindent\Large{\underline{\textbf{Exercise 5}}}
\normalsize{}\\
\noindent Let $\beta > 1$, solve in $\mathcal{C}^1(\R) \cap \mathcal{L}^1(\R)$ the following integral equation:
$$\int_{-\infty}^{+\infty} y(u)\, e^{-\beta(t-u)^2}\, du = e^{-t^2}, \qquad \text{for any } t \in \R.$$


% \input{Chapters/Chapter 8.tex}
% \input{Chapters/Series 8.tex}

\end{document}
